\documentclass[aspectratio=169]{beamer}
\usetheme{Madrid}
\usecolortheme{default}
\usefonttheme{professionalfonts}

\usepackage[spanish]{babel}
\usepackage[utf8]{inputenc}
\usepackage{booktabs}
\usepackage{amsmath}
\usepackage{xcolor}
\usepackage{tikz}

% Colors
\definecolor{aspblue}{RGB}{30, 90, 160}
\definecolor{aspgray}{RGB}{90, 90, 90}
\definecolor{aspgreen}{RGB}{34, 139, 34}
\definecolor{aspred}{RGB}{180, 30, 30}

\setbeamercolor{frametitle}{bg=aspblue, fg=white}
\setbeamercolor{title}{fg=white}
\setbeamercolor{author}{fg=white}
\setbeamercolor{date}{fg=white}
\setbeamercolor{structure}{fg=aspblue}
\setbeamercolor{palette primary}{bg=aspblue}
\setbeamercolor{palette secondary}{bg=aspblue!80}
\setbeamercolor{palette tertiary}{bg=aspblue!60}

\setbeamertemplate{navigation symbols}{}
\setbeamertemplate{footline}[frame number]

\title{Construcción del Modelo de Predicción de aspEnergy}
\subtitle{Regresión Polinomial con Lasso}
\author{}
\date{}

\begin{document}

% ── Portada ────────────────────────────────────────────────────────────────
\begin{frame}
  \titlepage
\end{frame}

% ── 1. Selección de características ───────────────────────────────────────
\begin{frame}{Punto de partida: las características}

  \textbf{23 campos numéricos de aspAudioFeatures disponibles.}

  \bigskip
  \begin{columns}
    \column{0.48\textwidth}
      \textbf{Incluidas (20):}
      \small
      \begin{itemize}
        \item Mood: danceability, acoustic, aggressive, happy, party, relaxed, sad
        \item Timbre: gender, voice, instrumental
        \item Estilo: electronic, approachability, engagement
        \item Físicas: bpm, loudness, intensity
        \item Pitch: avg, max, median, std
      \end{itemize}

    \column{0.48\textwidth}
      \textbf{\color{aspred}Excluidas (3):}
      \small
      \begin{itemize}
        \item \texttt{key} — valor de texto (ej. ``Eb'')
        \item \texttt{chord} — valor de texto (ej. ``C\#'')
        \item \texttt{mode} — valor de texto
      \end{itemize}
      \bigskip
      \begin{block}{\small ¿Por qué importa?}
        \small Incluirlas en el modelo anterior causó coerción silenciosa a~0, inflando el R²~=~0.99 en entrenamiento pero colapsando a R²~=~0.548 fuera de muestra.
      \end{block}
  \end{columns}
\end{frame}

% ── 2. El enfoque: Polinomial + Lasso ─────────────────────────────────────
\begin{frame}{El enfoque: Expansión Polinomial + Lasso}

  \begin{block}{Motivación}
    La energía perceptual no es lineal: una canción con alto \textit{aggressive} \textbf{y} alto \textit{relaxed} simultáneamente es distinta de una con solo uno de esos rasgos.
  \end{block}

  \bigskip
  \begin{enumerate}
    \item \textbf{Expansión grado-2:} generar todos los términos cuadráticos y productos cruzados
    \[
      20 \text{ features} \;\longrightarrow\; 230 \text{ términos posibles}
    \]
    \item \textbf{Estandarización:} \texttt{StandardScaler} — media 0, desviación estándar 1
    \item \textbf{Lasso con validación cruzada:} \texttt{LassoCV} selecciona automáticamente el parámetro de regularización $\alpha$ y elimina los términos irrelevantes
    \[
      230 \text{ términos} \;\xrightarrow{\text{Lasso}}\; \mathbf{48 \text{ términos no nulos}}
    \]
  \end{enumerate}

\end{frame}

% ── 3. CV y Early Stopping ────────────────────────────────────────────────
\begin{frame}{Validación cruzada y Early Stopping}

  \textbf{¿Hasta qué punto vale la pena añadir complejidad?}

  \bigskip
  \begin{table}
    \centering
    \begin{tabular}{lcccc}
      \toprule
      \textbf{Configuración} & \textbf{R² (5-fold CV)} & \textbf{MAE} & \textbf{Ganancia} \\
      \midrule
      Grado-1 lineal (20 feat.)       & 0.9187 & 0.051 & — \\
      Grado-2 polinomial (20 feat.)   & 0.9519 & 0.040 & +0.033 \\
      Grado-2 + feat. ingenierizadas  & 0.9526 & 0.039 & \textcolor{aspred}{+0.0007} \\
      \bottomrule
    \end{tabular}
  \end{table}

  \bigskip
  \begin{alertblock}{Early Stopping}
    La ganancia al añadir \textit{log(bpm)}, ratios de mood y variables compuestas fue de solo \textbf{+0.0007} en R² — por debajo del umbral de 0.003. El proceso se detuvo en grado-2.
  \end{alertblock}

  \small
  \textit{Los 5 folds mostraron R² estable (desviación estándar = 0.0004), confirmando que el modelo no está sobreajustado.}

\end{frame}

% ── 4. La fórmula ─────────────────────────────────────────────────────────
\begin{frame}{La fórmula resultante}

  \[
    \widehat{\text{aspEnergy}} = \text{clip}\!\left(\,1.014 + \sum_{i} c_i \cdot t_i,\; 0,\; 1\right)
  \]

  \bigskip
  \begin{columns}
    \column{0.48\textwidth}
      \textbf{\color{aspgreen}Términos positivos (aumentan energía):}
      \small
      \begin{itemize}
        \item happy $\times$ sad \hfill (+0.367)
        \item danceability $\times$ relaxed \hfill (+0.284)
        \item aggressive $\times$ happy \hfill (+0.204)
        \item aggressive $\times$ instrumental \hfill (+0.139)
      \end{itemize}

    \column{0.48\textwidth}
      \textbf{\color{aspred}Términos negativos (reducen energía):}
      \small
      \begin{itemize}
        \item aggressive $\times$ relaxed \hfill (−0.359)
        \item relaxed \hfill (−0.291)
        \item party $\times$ sad \hfill (−0.270)
        \item relaxed $\times$ voice \hfill (−0.159)
      \end{itemize}
  \end{columns}

  \bigskip
  \begin{block}{Interpretación}
    La energía emerge del \textbf{contraste emocional}, no de una sola dimensión. Una canción agresiva-y-relajada a la vez \textit{cancela} energía; una happy-y-sad simultánea la \textit{amplifica}.
  \end{block}

\end{frame}

% ── 5. Resultados de validación ───────────────────────────────────────────
\begin{frame}{Resultados de validación}

  \begin{columns}
    \column{0.48\textwidth}
      \begin{block}{Validación interna (137k canciones)}
        \begin{itemize}
          \item \textbf{R²} = 0.9525
          \item \textbf{MAE} = 0.039
          \item \textbf{Sesgo} $\approx$ 0 (±0.0004)
          \item Estable en los 5 folds
        \end{itemize}
      \end{block}

      \bigskip
      \begin{block}{18 canciones fuera de base (zero-shot)}
        \begin{itemize}
          \item \textbf{R²} = 0.9564
          \item \textbf{MAE} = 0.025
          \item 17 de 18 canciones con error $<$ 0.05
        \end{itemize}
      \end{block}

    \column{0.48\textwidth}
      \begin{exampleblock}{¿Por qué el OOS es mejor que el CV?}
        El R² fuera de muestra (0.9564) es ligeramente \textit{mayor} que el CV interno (0.9525) — señal de que el modelo no está sobreajustado al conjunto de entrenamiento.
      \end{exampleblock}

      \bigskip
      \small
      \textbf{Prueba extra — 10 canciones de Naruto (J-pop/J-rock):}
      \begin{itemize}
        \item R² = 0.856, Pearson = 0.948
        \item El \textit{ranking} es casi perfecto; los valores absolutos se sobreestiman $\sim$0.04--0.11 (género subrepresentado en entrenamiento)
      \end{itemize}
  \end{columns}

\end{frame}

% ── 6. Lecciones aprendidas ───────────────────────────────────────────────
\begin{frame}{Lecciones del proceso}

  \begin{enumerate}
    \setlength{\itemsep}{10pt}

    \item \textbf{La calidad de datos fue el mayor riesgo.}\\
    \small Incluir \texttt{key} y \texttt{chord} (texto) infló el R² de entrenamiento a 0.99 pero destruyó la generalización (R²~=~0.548 OOS). Eliminarlos lo llevó a 0.9564.

    \medskip
    \item \textbf{El grado-2 es el punto óptimo de complejidad.}\\
    \small La ingeniería de características (log, ratios, variables compuestas) aportó solo +0.0007 en R². No valió la pena sacrificar interpretabilidad.

    \medskip
    \item \textbf{La validación cruzada es imprescindible.}\\
    \small Un único split train/test puede dar estimaciones optimistas. Los 5 folds confirmaron R² estable y eliminaron la duda de sobreajuste.

    \medskip
    \item \textbf{Una fórmula aritmética puede competir con un modelo de producción.}\\
    \small 48 términos, completamente transparentes e inspeccionables, alcanzan la misma precisión que el pipeline de ML de producción.

  \end{enumerate}

\end{frame}

\end{document}
